% Shared middle-defense body; the final thesis and experiment assets are read-only inputs.
\chapter{状态估计、控制与平台验证进展}\label{ch:middle-systems}

磁观测得到的横向中心、相对方向和距离信息，需要与载体运动状态相结合，才能转化为可执行的巡检目标。本阶段围绕这一衔接关系开展三类工作：在状态估计层处理观测退化与异常值，在决策控制层研究置信信息如何改变运动输出，在实物子链上核验磁采集、目标计算平台和执行脑通信。各类实验分别固定输入、比较条件和统计单位，形成从算法机制到实际接口的证据。以下先说明必要方法及其代表结果，再归纳原 AUV 平台的恢复状态。

\section{融合状态估计与异常观测处理}\label{sec:middle-estimation}

\subsection{状态、观测与不确定性的统一表达}

水下巡检同时需要载体自身状态与相对电缆几何。惯性器件提供连续运动递推，多普勒测速仪（Doppler velocity log，DVL）和深度计分别约束速度与垂向位置；有效的声学、磁学观测进一步提供目标相对位置或方向。不同观测的更新频率、延迟和失效方式不同，因此本研究采用误差状态扩展卡尔曼滤波（error-state extended Kalman filter，ES-EKF）作为统一融合骨架\cite{sola2017quaternion,madyastha2011eskf}。本阶段的工作重点是观测进入滤波器的条件、分源质量处理和不确定性传递，四元数误差状态及标准卡尔曼更新沿用既有理论。

名义状态写为
\begin{equation}
 \boldsymbol{x}=\left[\boldsymbol{p}^{n},\boldsymbol{v}^{n},
 \boldsymbol{q}_{b}^{n},\boldsymbol{b}_{a},\boldsymbol{b}_{g}\right],
 \qquad
 \delta\boldsymbol{x}=\left[\delta\boldsymbol{p},\delta\boldsymbol{v},
 \delta\boldsymbol{\theta},\delta\boldsymbol{b}_{a},\delta\boldsymbol{b}_{g}\right]^{\mathsf T}.
\end{equation}
其中，$n$ 为统一导航坐标系，$b$ 为载体坐标系，$\boldsymbol{p}^{n}$、$\boldsymbol{v}^{n}$ 分别为位置和速度，单位四元数 $\boldsymbol{q}_{b}^{n}$ 描述姿态，$\boldsymbol{b}_{a}$、$\boldsymbol{b}_{g}$ 为惯性器件偏置。名义状态含 16 个标量，受四元数单位约束而具有 15 个自由度；局部姿态误差由三维旋转矢量表示。预测步由惯性输入递推状态和协方差，观测到达后根据相应模型更新，再将误差注入名义状态。通过这一结构，姿态转动、传感器安装角与磁传感器杆臂可在统一坐标关系下处理。

对来源 $s$ 的观测，定义新息 $\boldsymbol{y}_{s}=\boldsymbol{z}_{s}-h_s(\hat{\boldsymbol{x}}^{-})$，新息协方差为 $\boldsymbol{S}_{s}=\boldsymbol{H}_{s}\boldsymbol{P}^{-}\boldsymbol{H}_{s}^{\mathsf T}+\boldsymbol{R}_{s}$。这里 $\boldsymbol{P}$ 描述状态不确定性，$\boldsymbol{R}_{s}$ 描述该来源的观测噪声。标准归一化新息平方（normalized innovation squared，NIS）及其按维数归一化形式为
\begin{equation}
 \epsilon_s=\boldsymbol{y}_{s}^{\mathsf T}\boldsymbol{S}_{s}^{-1}\boldsymbol{y}_{s},
 \qquad \bar\epsilon_s=\epsilon_s/m_s,
\end{equation}
其中 $m_s$ 为观测维数。标准 NIS 用于对应自由度的卡方区间检验，$\bar\epsilon_s$ 用于跨来源比较\cite{barshalom2001estimation}。深度是一维观测，DVL 是三维观测，将二者未经归一化的新息统计混入一个窗口，会使一种传感器的异常改变另一种传感器的权重。

本阶段采用分源窗口调节持续失配，并用 Huber 型即时权重抑制单帧脉冲。令 $r_s=\sqrt{\epsilon_s/m_s}$，当 $r_s$ 超过阈值 $\delta$ 时取 $w_s=\delta/r_s$，否则取 $w_s=1$，将当前观测协方差进一步调整为 $\boldsymbol{R}_{s,\mathrm{robust}}=\boldsymbol{R}_{s,\mathrm{adaptive}}/w_s$。前者处理一段时间内噪声水平的变化，后者限制单次异常对状态修正的影响。状态及其协方差同时提供给任务层和控制层；由协方差、观测质量映射得到的置信分数是运行决策量，其数值不直接解释为经过校准的真实成功概率。

\subsection{公平对照与鲁棒更新的代表结果}

首先在八类传感退化场景、每类三个种子的 24 个仿真数据包上，统一三种估计器的初始位置、速度、姿态与坐标变换，比较纯航位推算、标准 EKF 和 ES-EKF。三者水平 RMSE 的逐运行均值分别为 3.197、3.196 和 3.200 m，跨运行标准差分别为 0.695、0.692 和 0.692 m，均值加减一倍标准差区间高度重叠。该组输入仅含 DVL 速度和深度观测，缺少绝对水平位置约束，因此滤波结构本身没有消除速度积分漂移。纯航位推算较小的深度误差来自直接透传深度读数，也不构成其整体优于融合估计的证据。这组公平比较明确了本阶段的评价重点：连续输出、异常抑制和协方差诊断分别评价，水平定位改善还需由有效声磁几何约束提供信息。

针对异常抑制，另用受控低噪声轨迹，DVL 与深度分别以 5 Hz、10 Hz 更新，名义噪声标准差为 0.005 m/s 和 0.01 m；其绝对误差不与前述八场景矩阵横向比较。各策略采用同一运动轨迹、逐样本相同的噪声以及五个随机种子，设置名义、脉冲异常、脉冲叠加 100 ms 延迟三类场景。每次仿真持续 120 s，六种策略共形成 90 次离线滤波对照。表~\ref{tab:middle-filter-results} 中每格为五次运行三维 RMSE 的算术均值，统计单位是一次完整滤波运行，而非将每个时间点当作独立重复。

\begin{table}[htbp]
\centering
\caption{同输入滤波策略对照的三维 RMSE（m）}\label{tab:middle-filter-results}
\begin{tabular}{@{}lrrr@{}}
\toprule
策略 & 名义 & 脉冲异常 & 延迟异常 \\
\midrule
标准 EKF & 0.0214 & 1.1677 & 1.1719 \\
ES-EKF 固定 $R$ & 0.0251 & 1.1722 & 1.1763 \\
全局自适应 $R$ & 0.0251 & 1.6344 & 1.6398 \\
分源自适应 $R$ & 0.0251 & 1.2384 & 1.2304 \\
分源 $R$ 加 Huber & 0.0251 & 0.2440 & 0.2463 \\
再加时间戳噪声膨胀 & 0.0251 & 0.2440 & 0.2927 \\
\bottomrule
\end{tabular}
\end{table}

结果显示，名义条件下各 ES-EKF 变体保持相同结果，新增异常处理没有改变基线行为。脉冲条件下，历史全局自适应反而扩大误差；分源调节避免不同来源的统计混用，但仍不足以消除瞬时异常。加入 Huber 后，两类异常的 RMSE 较固定 $R$ 分别降低约 79\%，说明该组收益主要来自即时鲁棒降权。这一结论还需要与不确定性一致性分开：延迟异常下 Huber 方案的 DVL NIS/自由度均值仍为 2.768，95\% 卡方带覆盖率为 0.176。低误差并不意味着协方差已经正确标定。

时间戳处理的对照进一步给出设计约束。仅按观测年龄增大量测噪声，将延迟异常 RMSE 从 0.2463 m 增至 0.2927 m，说明仅按年龄降权未能弥补该组时序错位。因此，当前延迟感知加权是一项观测使用策略，不能替代历史状态更新。后续时序改进需要连接设备采样、传输到达和融合消费的同次记录，并在相同输入下单独检验；这一需求也与下文真实 PC104 时间基准核验相衔接。

\section{置信信息参与控制及其性能权衡}\label{sec:middle-control}

\subsection{控制调制与任务授权的职责划分}

状态估计存在不确定性时，控制器仍需决定是否继续跟踪、采用多大速度以及如何约束指令变化。本研究在模型预测控制（model predictive control，MPC）的有限时域优化框架\cite{mayne2000constrained}上，将感知置信度引入代价调度，形成不确定性感知 MPC（uncertainty-aware MPC，UA-MPC）。对置信度 $c\in[0,1]$，跟踪权重与控制代价缩放采用
\begin{equation}
\begin{aligned}
 w_e(c)&=w_{e,0}\bigl[1+s_e(1-c)^{\alpha}\bigr],\\
 s_u(c)&=\gamma_u+\frac{1-\gamma_u}{1+\exp[-k(c-c_0)]}.
\end{aligned}
\end{equation}
其中 $s_e$、$\alpha$ 控制低置信度下跟踪权重的变化，$\gamma_u$、$k$、$c_0$ 控制努力惩罚的平滑过渡。低置信度时提高对既有参考的约束，同时降低必要纠偏动作的代价，指令变化由速率与带宽约束限制。这是一种可解释的控制调制，并不等同于对未知环境建立了概率安全保证。

控制调制与任务授权使用不同的判据。行为树根据观测是否有效、目标是否明确和系统健康状态，决定保持、搜索或跟踪；控制器只在相应授权下消费目标，并由独立仲裁环节处理失联、故障和陈旧命令。声学来源用于目标跟踪质量，磁检测分数用于表征异常存在性，两者分别保留，以免把强磁峰自动变成指定目标的跟踪许可。正常任务切换和执行安全由不同层负责，使拒绝某次估计、降低运动强度以及撤销自治能够分别记录和分析。

\subsection{六类代理场景的成对比较}

代表性控制矩阵采用六类代理电缆场景，比较固定权重与保守分源 UA 两种模式，每组包含三个随机种子、每次持续 60 s，共 36 次运行。场景覆盖连续 S 弯、发卡掉头、坡面穿越、埋设间断、横流冲击和综合极端；相同场景种子下的两种模式共享输入设置。全部运行完成并生成有效记录，求解回退为零。这说明实验链与求解过程可持续运行，任务精度和安全表现另由对应指标判断。

\begin{table}[htbp]
\centering
\caption{六类代理场景中 UA 控制的收益与代价}\label{tab:middle-control-results}
\begin{tabular}{@{}lrrr@{}}
\toprule
场景 & 固定权重横偏 RMSE/m & UA 横偏 RMSE/m & 控制变化率变化 \\
\midrule
连续 S 弯 & 5.898 & 6.002 & $-24.5\%$ \\
发卡掉头 & 6.017 & 6.013 & $-89.8\%$ \\
坡面穿越 & 4.140 & 4.127 & $+15.9\%$ \\
埋设间断 & 4.136 & 4.053 & $+2.3\%$ \\
横流冲击 & 15.052 & 15.484 & $-52.0\%$ \\
综合极端 & 12.349 & 12.356 & $-84.4\%$ \\
\bottomrule
\end{tabular}
\end{table}

表~\ref{tab:middle-control-results} 给出逐场景三次运行均值。控制变化率按相邻有效时刻的控制向量差分除以时间间隔，取向量范数后计算 RMS；该控制向量含角度与推力分量，数值依赖固定的指令编码尺度，仅用于同契约成对比较，不解释为物理能耗。再对六类场景等权平均，固定权重与 UA 的控制变化率 RMS 分别为 0.247 和 0.085，横向 RMSE 则由 7.932 m 变为 8.006 m。发卡、横流和综合极端中的指令波动降低，而坡面与埋设间断没有同方向改善。因此，可保留的贡献是特定退化场景下的控制平滑性及分源调度能力，不能概括为跟踪精度普遍提升。

综合极端还保留了影响主结论的负结果：固定权重和 UA 的近底安全违规样本比例分别约为 0.290 和 0.287，UA 未进入跟踪授权态。该组结果意味着保守调度能够降低指令变化，却未同时完成复杂场景的安全巡检。评价时应分别报告运行完成、求解可解、任务授权、路径精度与安全违规，避免以单一“成功率”掩盖职责不同的指标。相应的控制修正与水中验证作为后续独立工作包，不通过修改统计窗口将其处理为正例。

\section{弱磁测量链的实测基础}\label{sec:middle-measurement}

\subsection{TMR 模组与 ADC 噪声评价}

仿真磁观测的噪声参数，需要由实际采集链提供量级依据。本阶段采用 TMR8637 三轴模组与 SK2301 数采链。编号 6\# 模组的随附测试报告给出三轴灵敏度分别为 20.02、19.98 和 19.96 mV/$\mu$T。这些参数作为电压到磁场的转换依据；传感器本体报告与本研究实施的 ADC 短接测试分开使用，以区分器件先验和集成后的测量证据。

ADC 短接实验将 CH4 差分输入短接，在各采样率下分别采集 10 s，排除目标磁场和传感器输出对这一子链的影响。有效位数沿用原实验的噪声口径
\begin{equation}
 \mathrm{ENOB}_{\mathrm{noise}}
 =\log_2\!\left(\frac{V_{\mathrm{span}}}{\sqrt{12}\,\sigma_V}\right),
\end{equation}
其中满量程跨度 $V_{\mathrm{span}}=20$ V，$\sigma_V$ 为去均值电压的 RMS。直接以 2 kHz 采样时，RMS 噪声为 47.974 $\mu$V，对应 16.877 bit；以 16 kHz 采样并采用过采样比 $\mathrm{OSR}=8$ 的非重叠块平均后，输出恢复为 2 kHz，RMS 降为 17.375 $\mu$V，ENOB 达到 18.342 bit。这里标称 24 bit 描述数字量化位宽，噪声 ENOB 才描述规定测试条件下的有效分辨能力。

\begin{figure}[htbp]
 \centering
 \includegraphics[width=0.79\linewidth]{../docs/thesis/figures/hardware/adc_enob/ch4_enob_vs_osr.pdf}
 \caption{SK2301 短接输入的噪声有效位数随过采样比变化；各硬件采样率分别采集 10 s，曲线用于评价数采子链。}\label{fig:middle-adc}
\end{figure}

图~\ref{fig:middle-adc} 中不同采样率下的曲线基本重合，表明块平均能够稳定降低宽带随机噪声。对磁目标采用相干积分还可进一步限定观测带宽：在 16 kHz 采集、OSR 为 8、1 s 的 45 Hz 正交解调条件下，按 20 mV/$\mu$T 名义灵敏度折算，十个窗口的矢量 RMS 为 0.0230 nT，I/Q 单分量的 $3\sigma$ 为 0.0514 nT。这支持数采子链达到约 0.05 nT 的窄带灵敏度量级；后一判据略高于 0.05 nT，不能写成所有噪声判据均已达标。完整测量链的准确度还需计入传感器、模拟前端、温漂、电源和几何标定。

短接测量与实际环境输入承担不同作用。已有汇总记录还报告了包含传感器、供电和接线的真实三轴背景记录，并将其作为仿真观测噪声源。该路线保留真实频谱和轴间相关性，有助于检验只采用独立高斯噪声时遗漏的时间结构。背景记录来自非屏蔽环境，其中工频分量属于环境扰动；这些记录用于回放压力与接口验证，不能将其频谱直接视作模组本征噪声。测量链因此形成器件先验、短接子链和实际背景三个互补层次。

\subsection{45 Hz 电缆与相对位姿联合实验}

为检验磁场观测与空间几何能否形成一致约束，本阶段以 45 Hz 电流激励实验室直导线，同步记录三轴磁电压和双 ArUco 相对位姿。ArUco 在此作为实验室几何测量工具，用于将移动传感器的观测与导线坐标对齐，并非水下任务的既定定位传感器。ADC 以 2 kHz 记录 106821 个时刻样本；按 1 s 锁相窗与 0.05 s 步长完成对齐，得到 777 个有效磁场--位姿配对点。窗口高度重叠，故 777 表示有效配对点数，不作为 777 次独立重复实验。

黄金电阻两端电压为 7.56 V RMS，电阻为 4.09 $\Omega$，对应电流为 1.8484 A RMS 和 2.6140 A 峰值。锁相输出采用正弦峰值口径，模型因此采用峰值电流。磁场由无限长直导线模型给出：以单位向量 $\boldsymbol{u}$ 表示导线轴向，以 $\boldsymbol{p}_0$ 表示轴线上一点，将传感器相对位移投影到导线法平面得到 $\boldsymbol{\rho}$，则
\begin{equation}
 \boldsymbol{B}(\boldsymbol{p})=\frac{\mu_0 I}{2\pi}
 \frac{\boldsymbol{u}\times\boldsymbol{\rho}}{\|\boldsymbol{\rho}\|^2},
 \qquad
 \boldsymbol{\rho}=(\boldsymbol{I}-\boldsymbol{u}\boldsymbol{u}^{\mathsf T})
 (\boldsymbol{p}-\boldsymbol{p}_0).
\end{equation}
拟合保留三轴同相和正交分量，以公共复尺度表达电流幅值与参考相位，联合估计直线的可辨识几何自由度。相比仅拟合非负磁场模量，复 I/Q 保留了方向符号和轴间相位关系，使传感器运动过程中的峰谷、换向共同约束导线位置。

\begin{figure}[htbp]
 \centering
 \includegraphics[width=\linewidth]{middle/figures/systems/complex_iq_fit_middle.pdf}
 \caption{实验室移动 TMR 的三轴 45 Hz 复 I/Q 与直导线模型拟合。左、右列分别为同相和正交分量；777 个配对点来自同一段连续记录。}\label{fig:middle-joint}
\end{figure}

图~\ref{fig:middle-joint} 表明，模型能解释三轴观测的主要峰谷与换向，但部分区段仍有幅值和相位残差。自由复尺度拟合的复数决定系数 $R^2=0.8792$，等效峰值电流为 2.8789 A；固定实测峰值电流为 2.6140 A 时，$R^2=0.8711$。固定电流后仍保持接近的解释水平，支持空间几何与磁场模型之间存在稳定联系；自由尺度较高的等效电流也提示通道增益、杆臂和模型近似不能仅靠一个拟合系数消除。

该实验完成的是移动磁传感器、同步几何与矢量磁模型的联合验证，为受控横切与安装外参分析提供了实际数据基础。其自由手持轨迹不同于平面、双侧完整的定高横切；在后续任务设计中，应把“存在可解释的空间变化”与“达到 FWHM 深度发布条件”分别评价，使同一段真实数据既能支持已完成的几何分析，也能用于检查不充分扫描的拒绝逻辑。

\section{目标平台计算与真实通信安全子链}\label{sec:middle-runtime}

\subsection{Jetson 求解预算及压力工况}

算法可求解还需与目标平台的更新时间预算相匹配。本阶段在 Jetson Orin NX 的 25 W、8 核模式下，对预测步数 $N=20$、离散步长 $\Delta t=0.2$ s 的 MPC 入口实施独立微基准。正式末轮稳态冷、热启动各运行 200 次；初始深度 8 m 到目标深度 3 m 的约束压力档各运行 50 次。wall-clock 耗时覆盖一次求解入口的实际执行，区别于优化器内部计时，图~\ref{fig:middle-jetson} 同时列出耗时分布与求解成功率。

\begin{figure}[htbp]
 \centering
 \includegraphics[width=\linewidth]{middle/figures/systems/jetson_mpc_latency_and_success_middle.pdf}
 \caption{Jetson Orin NX 独立 MPC 求解微基准：25 W、8 核，稳态冷/热启动各 200 次，压力档各 50 次。50 ms 为 20 Hz 更新预算，200 ms 为模型离散步长。}\label{fig:middle-jetson}
\end{figure}

稳态冷、热启动的耗时 p95 分别为 36.560 和 34.118 ms，全部求解收敛，在该状态附近低于 50 ms 的控制周期。约束压力下冷、热启动也全部收敛，但 p95 升至 122.342 和 119.795 ms，已超过这一预算。由此可见，求解成功与按时返回是不同要求；困难工况必须在执行链中配套截止时间管理、异步更新或回退策略，不能只凭优化器返回成功就继续消费迟到的结果。

同一提交、功耗与核数下的连续三轮记录，稳态热启动 p95 处于 33.948--35.464 ms，压力热启动为 119.795--122.230 ms，说明两类工作点之间的耗时差异能够短时重复。上述统计针对目标机单求解入口。与其独立完成的真实 PC104 通信实验采用另一宿主机中继拓扑，两者耗时不相加为整条磁采集--估计--决策--控制链的端到端延迟。目标平台的整栈调度及持续热负载由后续工作单独验收。

\subsection{PC104 固件回显与跨层故障传播}

沿用原 AUV 的 PC104/VxWorks 执行脑，可在保留既有底层接口的同时接入新的上层决策控制。实际通信核验以容器、macOS 宿主机双向中继和真实 PC104 组成链路，下行固定为零推进、零舵命令。中继在转发前反解协议并检查执行量，板端快照和上行反馈提供另两层核对。这样可以在执行机构不动作的条件下，检验报文解释、固件接收与故障状态传播是否一致。

300 s 固件回显记录以 10 Hz 发送 3000 帧下行，解析 4505 帧上行且无解析错误。对重复携带的回显按“下行帧号、固件接收 uptime”去重，得到 1001 个可配对接收事件。首次回显应用路径往返时间的 p95 为 313.965 ms，固件同一时钟域内接收到首次打包的 p95 为 16 ms。相邻唯一接收事件间隔均为 300 ms，反映当前固件接收任务约 3.33 Hz 的处理节拍；1001/3000 的比例是应用层观测覆盖率，不解释为物理网络丢包率。

这组时序结果直接影响接口选择：低频目标量下发与执行脑本地内环保持能够分担网络节拍约束，全执行量透传则必须考虑固件接收带宽。首次回显包含主机、中继、双向网络及固件调度，两个设备未共享绝对时钟，故不能由往返时间拆出严格单向物理时延。PC104 uptime 相对主机时钟的拟合速率差约为 $+962$ ppm，也提示长期融合需要估计时钟斜率，而不能长期沿用一次性的时间偏移锚定。

在正常通信基础上，将真实 PC104 故障字接入生产桥接、命令仲裁器与行为树，使用同一记录时间轴观测状态传递。DVL 失锁位 Bit13 使精确跟踪条件失效并切换搜索，但保留自治授权；系统通信故障 Bit5 和 Jetson 看门狗 Bit14 则撤销自治并进入待机。图~\ref{fig:middle-pc104} 同时展示故障位、仲裁状态、任务输出和零执行量，避免只看到某一层日志就认定安全链已经贯通。

\begin{figure}[htbp]
 \centering
 \includegraphics[width=\linewidth]{middle/figures/systems/pc104_cross_layer_fault_timeline_middle.pdf}
 \caption{真实 PC104--ROS2--行为树同步故障响应。横轴相对健康自主基线计时，响应延迟在正文按各故障阶段起点另计；所有推进及舵命令为零，竖向虚线标示显式重新授权。本图沿用已归档的矢量曲线。}\label{fig:middle-pc104}
\end{figure}

修正后的短时同步运行通过 18 项验收。计时从各故障阶段起点开始，到主机首次观测满足条件的状态为止：Bit13 阶段的仲裁、行为树响应分别为 65.097 和 349.100 ms；Bit5 阶段的拒绝授权、待机分别为 102.532 和 188.463 ms；暂停桥接下行后，首次观测看门狗故障及拒绝授权为 1071.706 ms，包含看门狗等待时间。故障清除后系统保持拒绝，收到显式零命令授权才恢复自治。这一恢复规则防止仅因故障位短时回落而重新释放执行指令。

另一次独立运行完成 1800 s 零执行器稳态，124841 个仲裁状态样本均处于有效授权，目标故障状态样本为零。该包稳态之后的旧恢复附加序列仍保留一项失败，故只能认定完整稳态子区间通过；它与修正后的 18 项故障运行不能拼成一次全部通过的长时故障实验。通信台架对电源门限及 DVL 状态采用了临时隔离条件，因此本阶段形成的是实际执行脑参与的通信和同步安全状态链证据，能源健康与真实反馈的验收在整艇恢复中另列。

\section{整艇恢复现状与下一验证入口}\label{sec:middle-platform-status}

本课题以原有 AUV 为研究平台。原艇结构、能源和执行机构的历史设计提供集成约束，作者新增工作主要在感知载荷、计算通信、估计控制及相应软件验证，不能将原厂设计与历史性能直接列为本阶段成果。从当前可核查材料看，磁采集和实验室几何联合记录、Jetson 求解器、真实 PC104 通信与零执行器故障传播已有证据；整艇能源、必要导航反馈、执行机构动态和水中运行则尚未形成相同层级的验收链。

下一工程阶段因此转向基础平台恢复与水中受控运行。新增计算平台和磁载荷进入实际执行链的顺序，将依据现场接口与准入条件确定；已有双脑软件设计为该决策提供接口基础。

对已完成子链，下一步应复用其明确输入输出，例如以固件回显核对真实接令，以故障位与仲裁同步记录核对撤权，以磁原始数据和位姿记录核对测量几何。对尚未验证的整艇环节，则保存供电、反馈、实际动作与退出过程的同次证据，再逐级增加运行自由度。这样既能利用现有工作降低联调成本，也能明确区分基础受控运行与海缆自主任务。具体依赖、验收条件和未满足时的替代路径在后续计划中集中安排。
